% Source-aligned English for printed p.131.
\SourcePage{131}

\LayerTarget{tohoku:def:additive-diagram-relations:p2}
\noindent $R_{ij}$, we shall say that $D$ is $R$-commutative if every $D(L)$,
$L\in R$, is zero. A pair $(S,R)=\Sigma$, consisting of a diagram scheme and a
set $R$ as above, will be called a \textit{diagram scheme with commutation
relations}. For every additive category $C$, we may then consider the
subcategory $C^\Sigma$ of $C^S$ consisting of the $R$-commutative diagrams.

\LayerTarget{tohoku:prop:1.6.1}
\noindent\textsc{Proposition 1.6.1.} --- \textit{Let $\Sigma$ be a diagram
scheme with commutation relations, and let $C$ be an additive category. Then
$C^\Sigma$ is an additive category. If $C$ is a category with infinite direct
products (respectively, infinite direct sums), then so is $C^\Sigma$.
Moreover, if $C$ satisfies one of axioms AB 1) through AB 6), or one of the
dual axioms AB 3*) through AB 6*), then so does $C^\Sigma$.}

\LayerTarget{tohoku:prop:1.6.1:proof}
Indeed, if $D,D'\in C^\Sigma$ and $u$ is a morphism from $D$ to $D'$, then its
kernel (respectively, cokernel, image, or coimage) is the diagram formed by the
kernels (respectively, \ldots) of the components $u_i$; the morphisms in this
diagram, corresponding to the arrows of the scheme, are deduced from those of
$D$ (respectively, $D'\ldots$) in the evident way by restriction
(respectively, passage to the quotient). Direct sums and direct products of
families of diagrams are interpreted analogously. The subobjects $D'$ of a
diagram $D$ are identified with families $(D'(i))$ of subobjects of the
$D(i)$ such that, for every arrow $\varphi$ with origin $i$ and endpoint $j$,
$D(\varphi)D'(i)\subset D'(j)$; then $D'(\varphi)$ is defined as the morphism
$D'(i)\longrightarrow D'(j)$ induced by $D(\varphi)$. Quotient objects of $D$
are determined dually.

\LayerTarget{tohoku:def:dual-diagram-scheme}
If $S$ is a diagram scheme, the scheme having the same vertices and the same
set of arrows as $S$, but with the origin and endpoint of every arrow
interchanged, is called the \textit{dual scheme} and is denoted by $S^\circ$.
If a set $R$ of commutation relations for $S$ is also given, we retain the
same set for $S^\circ$. With this convention, for an additive category $C$,
the dual category of $C^\Sigma$ is identified with
$(C^\circ)^{\Sigma^\circ}$.

\LayerTarget{tohoku:def:canonical-diagram-extension}
Let $C,C'$ be two additive categories and let $\Sigma$ be a diagram scheme
with commutation relations. For every functor $F$ from $C$ to $C'$, one
defines in the evident way the functor $F^\Sigma$ from $C^\Sigma$ to
$C'^\Sigma$, called the \textit{canonical extension of $F$ to diagrams}.
Formally, $F^\Sigma$ behaves functorially in the argument $F$; in particular,
a natural transformation $F\longrightarrow F'$ defines a natural
transformation $F^\Sigma\longrightarrow F'^\Sigma$. Finally, for a composite
functor one has

\[
  (GF)^\Sigma=G^\Sigma F^\Sigma,
\]

\noindent and the exactness properties of a functor are preserved by
extension to a class of diagrams.

\LayerTarget{tohoku:section:1.7}
\noindent\textbf{1.7. Examples of categories defined by diagram schemes.}

\LayerTarget{tohoku:example:1.7:a}
\noindent a) Take $I$ to consist of one element and the set of arrows to be
empty. The commutation relations are then of the form $n_i e=0$, and hence
may be reduced to a single relation $ne=0$. Thus $C^\Sigma$ is the subcategory
of $C$ consisting of the objects annihilated by the integer $n$. If $n=0$,
this is $C$ itself.

\LayerTarget{tohoku:example:1.7:b:p1}
\noindent b) Take an arbitrary $I$, with no arrows and no commutation
relations. Then $C^\Sigma$ is identified with the product category $C^I$. If
there are given

\TohokuContinuation{132}
